\begin{abstract}
We evaluate a production-deployed retrieval-augmented question-answering
system over the Bank of Korea Compensation Regulations on a 100-question
benchmark spanning ten categories and four large language model
families. The system fields four backends — a hyper-relational
knowledge graph (TypeDB v3), a property graph (Neo4j), a preprocessed
context retriever, and a closed-book baseline — alongside four
in-context surface forms in which the entire knowledge graph is dumped
into the prompt as TypeQL, Cypher, JSON-LD, or natural language.
Counterintuitively, the production knowledge-graph agents underperform
the in-context surface forms by 21--30 percentage points on multi-key
queries (paired McNemar $p<0.05$ in 5 of 8 comparisons), despite the
agents having strictly more data. Disentangling \emph{query generation}
from \emph{KG expressivity}, we find that natural-language summary of
the same KG facts is Pareto-optimal in the (accuracy, token cost)
frontier (88.75\% mean accuracy, 14K vs 32K characters), and that one
LLM family (Llama-4-Scout) exhibits a 21-percentage-point form
sensitivity that other families do not — directly tracing the bottleneck
to query-language familiarity rather than KG expressivity. We further
show that in-context KG dumps are paradoxically weaker than natural
language at \emph{closed-world boundary inference} (proving absence),
because the schema-driven structural clarity of the KG is stripped when
serialized as a raw dump. We argue that KG infrastructure should be
valued as the \emph{operational backbone} of the system rather than as
the LLM-facing surface form, and propose explicit boundary annotations
in KG-to-natural-language conversion pipelines as a practical remedy.
\end{abstract}
